The needs of Internet users have been growing rapidly, surpassing the capacity of Internet Service Providers (ISPs) to adequately meet them. Significant variations in demand at different times of the day have led ISPs to seek solutions based on approaches that provide greater resource flexibility and improve resource management.
Network slicing is a technique aimed at improving the Quality of Service (QoS) provided by ISPs by partitioning network resources and functions into virtual slices designed to meet the specific requirements of different users and applications. This work provides an introductory overview of the concept of network slicing and the approaches that, together, have enabled its evolution. Subsequently, different network slicing algorithms are presented and their performance is compared in order to obtain more accurate information about the characteristics of each approach and identify the scenarios in which each algorithm performs best.